% Part: sets-functions-relations
% Chapter: sets
% Section: basics

\documentclass[../../../include/open-logic-section]{subfiles}

\begin{document}

\olfileid[pt]{sfr}{set}{bas}
\olsection{Extensionalidade}

Um \emph{conjunto} é uma coleção de objetos, considerada como um único
objeto. Os objetos que compõem o conjunto chamam-se \emph{elementos} ou
\emph{membros} do conjunto. Se $x$ é !!a{element} de um conjunto~$A$,
escrevemos $x \in A$; caso contrário, escrevemos $x \notin A$. O conjunto que
não tem !!{element}s chama-se o \emph{conjunto vazio} e
denota-se~``$\emptyset$''.

\begin{explain}
Não importa como \emph{especificamos} o conjunto, ou como
\emph{ordenamos} os seus !!{element}s, ou mesmo quantas \emph{vezes}
contamos os seus !!{element}s. Só importa quais são os seus !!{element}s.
Codificamos isto no princípio seguinte.
\end{explain}

\begin{defn}[Extensionalidade]
  Se $A$ e $B$ são conjuntos, então $A = B$ se e só se
  todo !!{element} de~$A$ é também !!a{element} de~$B$, e vice-versa.
\end{defn}

A extensionalidade autoriza certa notação. Em geral, quando temos
objetos $a_{1}$, \dots, $a_{n}$, então $\{a_{1}, \dots, a_{n}\}$ é
\emph{o} conjunto cujos !!{element}s são $a_1, \ldots, a_n$. Enfatizamos
a palavra ``\emph{o}'', pois a extensionalidade nos diz que só pode haver
\emph{um} tal conjunto. De fato, a extensionalidade também autoriza o
seguinte:
  \[
    \{a, a, b\} = \{a, b\} = \{b,a\}.
  \] 
Isto traduz a ideia de que, quando consideramos conjuntos, não nos importa a
ordem dos seus !!{element}s, nem quantas vezes são especificados.

\begin{tagblock}{novice}
\begin{ex}
Sempre que temos vários objetos, podemos reuni-los num
conjunto. O conjunto dos irmãos de Richard, por exemplo, é um conjunto que
contém uma pessoa, e podemos escrevê-lo como $S=\{\textrm{Ruth}\}$.
O conjunto dos inteiros positivos menores que $4$ é $\{1, 2, 3\}$, mas também
se pode escrever como $\{3, 2, 1\}$ ou mesmo como $\{1, 2, 1, 2, 3\}$.
São todos o mesmo conjunto, pela extensionalidade. De fato, todo !!{element}
de $\{1, 2, 3\}$ é também !!a{element} de $\{3, 2, 1\}$ (e de $\{1,
2, 1, 2, 3\}$), e vice-versa.
\end{ex} 
\end{tagblock}

Com frequência especificaremos um conjunto por alguma propriedade que os seus
!!{element}s compartilham. Usaremos a seguinte notação abreviada para isso:
$\Setabs{x}{\phi(x)}$, onde $\phi(x)$ representa a propriedade
que~$x$ tem de satisfazer para ser contado entre os !!{element}s do
conjunto. 

\begin{tagblock}{novice}
\begin{ex}
No nosso exemplo, também poderíamos ter especificado $S$ como
\[
S = \Setabs{x}{x \text{ é irmão de Richard}}.
\]
\end{ex}
\end{tagblock}

\begin{tagblock}{math}
\begin{ex}
Diz-se que um número é \emph{perfeito} se e só se for igual à soma dos seus
divisores próprios (isto é, números que o dividem sem resto mas não são
iguais ao próprio número). Por exemplo, $6$ é perfeito porque os seus
divisores próprios são $1$, $2$, e~$3$, e $6 = 1 + 2 + 3$. De fato,
$6$ é o único inteiro positivo menor que $10$ que é perfeito. Assim,
usando a extensionalidade, podemos dizer:
\[
	\{6\} = \Setabs{x}{x\text{ é perfeito e }0 \leq x \leq 10}
\]
Lemos a notação à direita como ``o conjunto dos $x$ tais que $x$
é perfeito e $0 \leq x \leq 10$''. A identidade aqui confirma que,
quando consideramos conjuntos, não nos importa como são especificados.
E, mais em geral, a extensionalidade garante que há sempre
apenas um conjunto dos $x$ tais que $\phi(x)$.
Logo, a extensionalidade justifica chamar a
$\Setabs{x}{\phi(x)}$ \emph{o} conjunto dos $x$ tais que~$\phi(x)$.
\end{ex}
\end{tagblock}

A extensionalidade nos dá um modo de mostrar que conjuntos são idênticos: para
mostrar que $A = B$, mostre que sempre que $x \in A$ então também $x \in B$,
e sempre que $y \in B$ então também $y \in A$.

\begin{prob}
Prove que há no máximo um conjunto vazio, isto é, mostre que se $A$ e $B$
são conjuntos sem !!{element}s, então $A = B$.
\end{prob}

\end{document}
